% =================================================================
% Kapitel: Fazit und Ausblick
% =================================================================
\label{chap:fazit}

\section{Fazit}

Ziel der Arbeit war die Entwicklung zweier Reglerkonzepte für den Antrieb eines
Greifers: einer Positionsregelung für die Öffnungsweite und einer Regelung der Greifkraft.
Beide Aufgaben wurden gelöst, eine davon auf anderem Weg als ursprünglich vorgesehen.

Die Positionsregelung wurde wie geplant entworfen, simuliert und umgesetzt. Ihre
Auslegung folgt unmittelbar aus der Eigenschaft der Regelstrecke: Da der Treiber im
Geschwindigkeitsmodus betrieben wird, ist die Strecke ein reiner Integrator, und ein
Proportionalregler erreicht ohne Integralanteil stationäre Genauigkeit. Die erreichte
Abweichung von unter 0,05~mm unterschreitet die geforderten $\pm 0{,}5$~mm um den Faktor
zehn und entspricht der Breite der Totzone, die ihrerseits durch das Verhältnis von
Mikroschrittauflösung zu Encoderauflösung nach unten begrenzt ist. Die erreichte
Genauigkeit ist damit nicht durch die Reglerparametrierung, sondern durch die gewählte
Hardware bestimmt.

Für die Greifkraft war ein zweiter, geschlossener Regelkreis über die StallGuard-Funktion
des Treibers vorgesehen. Dieses Konzept wurde vollständig ausgelegt, simuliert und
implementiert, erwies sich am Prüfstand jedoch als nicht tragfähig: Die Messgröße zeigte
über den gesamten Kraftaufbau keine Lastabhängigkeit. Die Kraftaufgabe wurde daraufhin
über den bereits validierten Positionsregler und die bekannte Federkennlinie gelöst. Die
Zielkraft wird dabei in eine Zielposition umgerechnet und angefahren. Die erreichte
Abweichung liegt bei höchstens 1,6~\% und ist vollständig durch die Positioniergenauigkeit
erklärbar.

Diese Lösung ist formal keine Kraftregelung, sondern eine modellbasierte Vorsteuerung mit
geregelter Position. Sie erfordert die Kenntnis der Nachgiebigkeit des Objekts und kann
Änderungen der Kraft nach dem Greifvorgang nicht erkennen. Sie kommt dafür ohne
zusätzliche Sensorik aus und macht die geforderte Umschaltung zwischen Positions- und
Kraftbetrieb strukturell trivial, da beide Betriebsarten denselben Regler verwenden und
sich lediglich in der Herkunft des Sollwerts unterscheiden.

Tabelle \ref{tab:anforderungsabgleich} stellt die Anforderungen dem erreichten Stand
gegenüber.

% ---------------------------------------------------------------
\begin{table}[htbp]
    \centering
    \caption{Abgleich der Anforderungen mit dem erreichten Stand}
    \label{tab:anforderungsabgleich}
    \begin{tabular}{lp{8.5cm}}
        \hline
        Anforderung & Stand \\
        \hline
        Echtzeitfähigkeit
            & überwiegend erfüllt. Kein Drift, Ursache des Jitters bekannt und
              reproduzierbar; jeder fünfte Zyklus überschreitet die Periodendauer \\[0.3em]
        Positioniergenauigkeit
            & erfüllt. Gemessen unter 0,05~mm bei gefordert $\pm 0{,}5$~mm \\[0.3em]
        Regelbare Greifkraft
            & erfüllt. Abweichung höchstens 1,6~\%, stufenlos vorgebbar \\[0.3em]
        Überlast- und Kollisionsschutz
            & erfüllt \\[0.3em]
        Nahtlose Umschaltung
            & erfüllt \\[0.3em]
        Wiederholgenauigkeit
            & für die Kraftvorgabe über mehrere Läufe belegt; ein Langzeittest über eine
              größere Zahl von Greifzyklen wurde nicht durchgeführt \\[0.3em]
        Dokumentation
            & erfüllt im Rahmen der angepassten Aufgabenstellung \\
        \hline
    \end{tabular}
\end{table}
% ---------------------------------------------------------------

\section{Methodische Erkenntnis}

Über das technische Ergebnis hinaus hat die Arbeit einen Befund hervorgebracht, der
allgemeiner gilt als der konkrete Anwendungsfall.

Das Simulationsmodell des verworfenen Kraftregelkreises war strukturell korrekt. Die
Wandlungskette von der Position über Federweg und Kraft zum Motormoment, die
Vorzeichenumkehr bei der Fehlerbildung, die Quantisierung der Messgröße und die Kopplung
der Abtastung an das Vollschrittraster bildeten das reale Verhalten zutreffend ab. Die
Simulation sagte dennoch ein Verhalten voraus, das am realen Aufbau nicht eintrat.

Die Ursache lag nicht in der Struktur des Modells, sondern in zwei Zahlenwerten: dem
Ruhewert der Sensorkennlinie und ihrer Steigung. Beide waren Annahmen, die nie an der
Hardware überprüft wurden, weil zum Zeitpunkt der Modellbildung kein Prüfstand existierte.
Die Simulation konnte damit zeigen, dass die \emph{Regelstruktur} bei einem Messglied mit
den angenommenen Eigenschaften funktionieren würde -- nicht aber, ob das Messglied diese
Eigenschaften tatsächlich besitzt.

Die Aussagekraft eines Modells ist durch denjenigen seiner Parameter begrenzt, der am
schlechtesten abgesichert ist. Im vorliegenden Fall war das nicht der Regler, sondern der
Sensor. Die in \texttt{params.m} durchgeführte Kennzeichnung der Parameterherkunft nach
gemessenen, berechneten und angenommenen Werten hätte diesen Umstand bereits vor der
Inbetriebnahme sichtbar gemacht. Sie sollte nicht als dokumentarischer Aufwand, sondern
als Teil der Modellbildung verstanden werden.

\section{Ausblick}

Aus dem erreichten Stand ergeben sich mehrere Ansatzpunkte für eine Fortführung.

\subsection*{Automatische Bestimmung des Kontaktpunkts}

Der Kontaktpunkt $x_0$ wird derzeit manuell gesetzt und geht nach Tabelle
\ref{tab:fehlerquellen} unmittelbar in die Kraft ein. Er ist damit die kritische Größe des
Verfahrens. Eine automatische Erkennung ist über den Lastwinkel möglich, also über die
Differenz zwischen dem Schrittzähler des Treibers und der gemessenen Encoderposition.
Dieser liegt bei Leerfahrt praktisch bei null und springt bei Berührung auf einen
messbaren Wert. Die dafür erforderlichen Größen werden bereits erfasst.

\subsection*{Beseitigung des Abtastjitters}

Die in Abschnitt \ref{sec:ergebnis_echtzeit} beschriebene Störung der Zykluszeit geht
ausschließlich auf den blockierenden Zugriff auf \texttt{SG\_RESULT} zurück. Da diese
Größe für die umgesetzte Kraftvorgabe nicht mehr benötigt wird, ließe sich die Ursache
durch Entfernen des Zugriffs aus dem Reglertakt vollständig beheben. Alternativ würde eine
Erhöhung der Baudrate zum Treiber die Dauer des Zugriffs auf unter ein Drittel verkürzen.
Beide Maßnahmen sind mit geringem Aufwand umsetzbar und würden die Echtzeitanforderung
vollständig erfüllen.

\subsection*{Überprüfung am fertigen Greifer}

Die Verwerfung der Kraftregelung über StallGuard stützt sich auf Messungen bei rund 16~\%
Motorauslastung. Am fertigen Greifer, an dem beide Zahnstangen belastet werden, läge die
Auslastung bei etwa 38~\%. Eine Bestätigung des Befunds unter diesen Bedingungen steht aus.
Ebenso ist zu prüfen, ob die Kraftvorgabe über das Federmodell an den vorgesehenen
nachgiebigen Fingern die am Prüfstand erreichte Genauigkeit beibehält.

\subsection*{Weitere Punkte}

Ein Langzeittest über eine größere Zahl aufeinanderfolgender Greifzyklen wurde nicht
durchgeführt und wäre für eine Aussage zur Langzeitstabilität erforderlich. Darüber hinaus
wäre für empfindliche Objekte eine Reduktion der Zustellgeschwindigkeit nach dem
Kontaktpunkt sinnvoll, da die Kraft bei voller Geschwindigkeit mit rund 85~N/s aufgebaut
wird.